\section{Question and results}

We study the passage from homogeneous constraints to inhomogeneous perturbations on fixed compact spaces in source-free Lorentzian Einstein--Cartan (EC) theory with constant complete Nieh--Yan (NY) coupling. Elimination of the independent connection, the linear gauge quotient, and tangency to an exact vacuum family are distinct stages of the same variational problem.

We use the action, quotient and boundary conventions of Papers01--05 \cite{paper1,paper2,paper3,paper3ec,paper4,paper5}. Paper06 treated the full equations of common-scale families, the homogeneous volume--shape cone and its exact homogeneous realization \cite{paper6}. Here these results provide the starting point for spatially dependent fields, closed harmonic blocks and background-dependent integral conditions.

First, the regular homogeneous metric strata have 5/3/2/2 canonical pairs on $T^3$/Nil/Sol/$S^3$. The static torus and symmetric Taub points are treated separately. For nonzero flat-torus modes, direct reduction before and after Lorentz gauge fixing and connection elimination gives the same quadratic dynamics and reduced two-form.

Second, the four flat-torus Taub conditions are necessary and sufficient for local integrability under the hypotheses of the Arms--Marsden--Moncrief theorem \cite{ArmsMarsdenMoncrief1982}. We give explicit second-order corrections for a standing wave with homogeneous expansion and retain all thirteen generated wavevectors for two nonparallel waves. The theorem application and the finite-order construction have different logical roles.

Third, we reduce inhomogeneous constraints on selected time-dependent vacuum backgrounds of Nil, Sol and $S^3$. We treat the central-zero and central-nonzero Nil sectors, fiber-independent Sol modes and general nonzero dual orbits, and four low Wigner labels on $S^3$, with complete reconstruction for $j=3/2$. Finite-time examples include an independent Lorenz evolution on Sol with interior convergence comparisons.

Two propagating field species, the total canonical pairs in a representation block, and the homogeneous regular-stratum count are different quantities. For Nil and Taub we establish quadratic necessary KID conditions. On the fixed Sol background, absence of KIDs and the existing theorem give local integrability of all full linearized constraint data in the stated Sobolev class (Section 7 and Appendix G). All-time stability and a complete physical spectrum are not established. Closed forms, finite metric paths and spatial Laplacians remain separate geometric data in Appendix E.

Auxiliary connection elimination, the two gravitational polarizations and linearization stability belong to established theory \cite{Hehl1976,Pons2010,BrillDeser1973,ArmsMarsdenMoncrief1982,FischerMarsdenMoncrief1980}. The concrete content here is the explicit constraints, reductions, reconstructions and second-order corrections for the stated action and quotients, with their counting units and domains.
